% Part: first-order-logic
% Chapter: syntax-and-semantics
% Section: subformulas

\documentclass[../../../include/open-logic-section]{subfiles}

\begin{document}

% SEGMENT OLP-0155-S01
\olfileid{fol}{syn}{sbf}

\olsection{\printtoken{P}{subformula}}

\begin{explain}
கொடுக்கப்பட்ட ஒரு !!{formula} வை ``உருவாக்கும்'' !!{formula}s பற்றி பேசுவது
பல நேரங்களில் பயனுள்ளதாக இருக்கும். இவற்றை அதன்
\emph{!!{subformula}s} என்கிறோம். எந்த !!{formula} வும் தன்னையே !!a{subformula}
ஆகக் கொண்டுள்ளது; $!A$ இன் $!A$ அல்லாத துணைச்சூத்திரம்
\emph{சரியான !!{subformula}}.
\end{explain}

\begin{defn}[உடனடி !!^{subformula}]
$!A$ ஒரு !!a{formula} ஆக இருந்தால், $!A$ இன்
\emph{உடனடி !!{subformula}s} பின்வருமாறு தொகுத்தறிதல் வழியில்
வரையறுக்கப்படுகின்றன:
\begin{enumerate}
\item அணு !!{formula}s க்கு உடனடி !!{subformula}s எதுவும் இல்லை.

\tagitem{prvNot}{\indcase{!A}{\lnot !B}{$\indfrm$ இன் ஒரே உடனடி
    !!{subformula}~$!B$.}}{}

\item \indcase{!A}{(!B \ast !C)}{$\indfrm$ இன் உடனடி !!{subformula}s
  $!B$, $!C$ (இங்கு $\ast$ என்பது ஏதேனும் ஒரு இரு-இடத் தொடுப்புக்குறி).}

\tagitem{prvAll}{\indcase{!A}{\lforall[x][!B]}{$\indfrm$ இன் ஒரே உடனடி
    !!{subformula}~$!B$.}}{}

\tagitem{prvEx}{\indcase{!A}{\lexists[x][!B]}{$\indfrm$ இன் ஒரே உடனடி
    !!{subformula}~$!B$.}}{}
\end{enumerate}
\end{defn}

\begin{defn}[சரியான !!^{subformula}]
$!A$ ஒரு !!a{formula} ஆக இருந்தால், $!A$ இன்
\emph{சரியான !!{subformula}s} பின்வருமாறு மறுநிகழ்வு வழியில்
வரையறுக்கப்படுகின்றன:
\begin{enumerate}
\item அணு !!{formula}s க்கு சரியான !!{subformula}s எதுவும் இல்லை.

\tagitem{prvNot}{\indcase{!A}{\lnot !B}{$\indfrm$ இன் சரியான !!{subformula}s
    $!B$ மற்றும் $!B$ இன் அனைத்து சரியான !!{subformula}s உம் ஆகும்.}}{}

\item \indcase{!A}{(!B \ast !C)}{$\indfrm$ இன் சரியான !!{subformula}s
  $!B$, $!C$, மேலும் $!B$ மற்றும்~$!C$ இன் அனைத்து சரியான
  !!{subformula}s உம் ஆகும்.}

\tagitem{prvAll}{\indcase{!A}{\lforall[x][!B]}{$\indfrm$ இன் சரியான
    !!{subformula}s~$!B$ மற்றும்~$!B$ இன் அனைத்து சரியான !!{subformula}s உம்
    ஆகும்.}}{}

\tagitem{prvEx}{\indcase{!A}{\lexists[x][!B]}{$\indfrm$ இன் சரியான
    !!{subformula}s~$!B$ மற்றும்~$!B$ இன் அனைத்து சரியான !!{subformula}s உம்
    ஆகும்.}}{}
\end{enumerate}
\end{defn}

\begin{defn}[!!^{subformula}]
$!A$ இன் !!{subformula}s என்பது $!A$ தானும் அதன் அனைத்து சரியான
!!{subformula}s உம் சேர்ந்த தொகுப்பு.
\end{defn}

\begin{explain}
உடனடி !!{subformula}s மற்றும் சரியான !!{subformula}s ஆகியவற்றை நாம்
வரையறுத்த விதத்தில் உள்ள நுண்ணிய வேறுபாட்டைக் கவனியுங்கள். முதல் நிலையில்,
ஒரு சூத்திரம்~$!A$ இன் உடனடி !!{subformula}s ஐ நேரடியாக வரையறுத்தோம்;
$!A$ இன் சாத்தியமான ஒவ்வொரு வடிவிற்கும் இது நிகழ்வுகளால் அளிக்கப்பட்ட வெளிப்படையான
வரையறை; நிகழ்வுகள் !!{formula}s தொகுப்பின் தொகுத்தறிதல் வரையறையை
பின்பற்றுகின்றன. இரண்டாவது நிலையில், எல்லா !!{formula}s தொகுப்பு
வரையறுக்கப்படும் முறையையே பின்பற்றுகிறோம்; ஆனால் ஒவ்வொரு நிகழ்விலும்
சிறிய !!{formula}s $!B$, $!C$ தாமாக இருப்பதுடன் இந்த !!{formula}s இன்
சரியான !!{subformula}s ஐயும் சேர்க்கிறோம். இதனால் வரையறை
\emph{மறுநிகழ்வு} ஆகிறது. பொதுவாக, தொகுத்தறிதல் வழியில் வரையறுக்கப்பட்ட
ஒரு தொகுப்பில் (இங்கு !!{formula}s) ஒரு சார்பை வரையறுக்கும் போது, அந்த
சார்பின் வரையறை நிகழ்வுகளே அதே சார்பைப் பயன்படுத்தினால் அது மறுநிகழ்வு
வரையறை. இருப்பினும் நன்கு வரையறுக்கப்பட, நாம் வரையறுக்கும் வாதத்திற்கு
``முன்'' வரும் வாதங்களுக்கான சார்பு மதிப்புகளை மட்டுமே பயன்படுத்துகிறோம்
என்பதை உறுதிசெய்ய வேண்டும். இங்கு $(!B \ast !C)$ க்கான ``சரியான
!!{subformula}'' ஐ வரையறுக்கும் போது ``முந்தைய'' !!{formula}s $!B$, $!C$
இன் சரியான !!{subformula}s ஐ மட்டுமே பயன்படுத்துகிறோம்.
\end{explain}

\begin{prop}
\ollabel{prop:subfrm-trans}
$!B$ என்பது $!A$ இன் ஒரு subformula என்றும், $!C$ என்பது $!B$ இன் ஒரு
subformula என்றும் வைத்துக்கொள்வோம். அப்போது $!C$ என்பது $!A$ இன் ஒரு
subformula. வேறு சொற்களில், subformula தொடர்பு கடப்புநிலையானது.
\end{prop}

\begin{prob}
\olref[fol][syn][sbf]{prop:subfrm-trans} ஐ நிறுவுக.
\end{prob}

\begin{prop}
\ollabel{prop:count-subfrms}
$!A$ என்பது $n$ தொடுப்புக்குறிகளும் அளவையடைகளும் கொண்ட ஒரு சூத்திரம் எனில்,
$!A$ இல் அதிகபட்சம் $2n+1$ துணைச்சூத்திரங்கள் உள்ளன.
\end{prop}

\begin{prob}
\olref[fol][syn][sbf]{prop:count-subfrms} ஐ நிறுவுக.
\end{prob}

\end{document}
